%============================================================
% CH1 · PATTERN 1 — EXERCISES, HINTS, ANSWERS
%============================================================
\subsection{Pattern 1: Exercises and Going Further}
\label{ssec:ch1-p1-ex}

\begin{patternbox}{Pattern 1 Summary Card}
\begin{tabularx}{\linewidth}{@{}>{\bfseries}l X@{}}
Recognise & the answer for a vertex (or a pair) depends only on the edges touching it \\
Core loop & for every edge $(u,v)$: \texttt{deg[u]++, deg[v]++} \;(directed: \texttt{out[u]++, in[v]++}) \\
Tools & special values ($0$, $1$, $n-1$); net score $\text{in}-\text{out}$; pairs $\deg u + \deg v - [uv\in E]$; $\sum_{E}(f(u)+f(v)) = \sum_V f(v)\deg(v)$ \\
Law & handshake: $\sum_v \deg(v) = 2|E|$, so the number of odd-degree vertices is even \\
Cost & $\bigO{|V|+|E|}$ time, $\bigO{|V|}$ memory, no adjacency list \\
Top traps & one endpoint only; labels $1..n$ in an array of size $n$; shared edge counted twice; degrees used for connectivity \\
\end{tabularx}
\end{patternbox}

\exlegend
\begin{enumerate}[leftmargin=*, label=\textbf{E1.\arabic*}, itemsep=6pt]
    \item \exeasy\ \pconcept \exprob{Find Center of Star Graph}{LeetCode 1791}{https://leetcode.com/problems/find-center-of-star-graph/}\ \lc{Easy}. An undirected star with $n \ge 3$ vertices (one center joined to all others, $n-1$ edges) is given as its edge list. Return the center.
    \item \exeasy\ \precent \cf{800}\ \exprob{Destroying Bridges}{Codeforces 1944A}{https://codeforces.com/problemset/problem/1944/A}. Islands $1..n$ form a complete graph. An adversary destroys at most $k$ bridges to minimise the number of islands reachable from island 1 (including island 1). Print that minimum, for $t \le 1000$ test cases with $n \le 100$.
    \item \exmedium\ \pconcept \exprob{Maximum Total Importance of Roads}{LeetCode 2285}{https://leetcode.com/problems/maximum-total-importance-of-roads/}\ \lc{Medium}. Assign the values $1..n$ (each once) to the $n$ cities; a road's importance is the sum of its two endpoint values. Maximise the total importance of all roads. $n, |E| \le 5\cdot10^4$.
    \item \exmedium\ \ptricky \exprob{Maximum Star Sum of a Graph}{LeetCode 2497}{https://leetcode.com/problems/maximum-star-sum-of-a-graph/}\ \lc{Medium}. Vertices have values in $[-10^4, 10^4]$. A star is a center plus some of its neighbours; its sum is the sum of the values of its vertices. Return the maximum sum of a star with at most $k$ edges. $n, |E| \le 10^5$.
    \item \exeasy\ \ptricky \exown{Is it a degree sequence?} Which of the following can be the list of degrees of a simple undirected graph? (a) $(3,2,2,2)$, \ (b) $(3,3,1,1)$, \ (c) $(3,3,2,2)$, \ (d) $(4,1,1,1,1)$.
    \item \exmedium\ \pconcept \exown{Counting short paths.} (a) In an undirected simple graph, count the paths with exactly two edges $u - v - w$ ($u \ne w$, a path and its reverse counted once). (b) In a directed graph without opposite arrow pairs, count the directed paths $a \to b \to c$. Both in $\bigO{|V|+|E|}$.
    \item \exmedium\ \pconcept \exown{A tree with prescribed degrees.} Given $n \ge 2$ and $d_1,\dots,d_n \ge 1$ with $\sum d_i = 2n-2$, output the edges of a tree in which vertex $i$ has degree $d_i$ (one always exists).
\end{enumerate}

\subsubsection*{Hints}
\begin{enumerate}[leftmargin=*, label=\textbf{E1.\arabic*}]
    \item The center appears in \emph{every} edge. How many edges do you need to look at?
    \item Island 1 has degree $n-1$. How many bridges must be destroyed to separate it from even one island?
    \item Rewrite the total as a sum over cities. Which city should get the value $n$?
    \item Each center only looks at its own neighbours. Which neighbours can only hurt?
    \item Handshake parity first; then think about what a vertex of degree $n-1$ forces.
    \item Fix the middle vertex $v$ / $b$. How many ways are there to choose the two edges at it?
    \item A vertex of degree 1 is a leaf. Attach leaves to the vertices that still need edges; when such a vertex is satisfied except for one edge, it becomes a leaf itself.
\end{enumerate}

\subsubsection*{Answers}

\begin{solutionbox}[E1.1 Find Center of Star Graph]
The degree solution (count, return the vertex of degree $n-1$) is $\bigO{n}$. But every edge contains the center, so the center is the common endpoint of the first two edges --- $\bigO{1}$:
\begin{lstlisting}[style=cpp]
#include <bits/stdc++.h>
using namespace std;

class Solution {
public:
    int findCenter(vector<vector<int>>& edges) {
        int a = edges[0][0], b = edges[0][1];
        return (a == edges[1][0] || a == edges[1][1]) ? a : b;   // the shared endpoint
    }
};
\end{lstlisting}
Example 1, edges $[1,2],[2,3],[4,2]$: the first two share vertex $2$ $\Rightarrow$ 2. Example 2, $[1,2],[5,1],\ldots$ $\Rightarrow$ 1. This only works because the input is promised to be a star: with a general graph, use the degrees.
\end{solutionbox}

\begin{solutionbox}[E1.2 Destroying Bridges]
\textbf{Answer:} $1$ if $k \ge n-1$, otherwise $n$.

If $k \ge n-1$, destroy all $n-1$ bridges at island 1 (its whole degree): island 1 is alone. If $k < n-1$, every island $v \ne 1$ is still reachable: the direct bridge $1{-}v$ and the $n-2$ two-step routes $1{-}w{-}v$ form $n-1$ routes that share no bridge, and $k < n-1$ destructions cannot hit all of them. So there is nothing in between: either the full degree of island 1 is destroyed, or nothing changes.

\begin{center}
\begin{tikzpicture}[x=1cm, y=1cm]
    \node[gvsrc] (o) at (0,0) {1};
    \foreach \a/\l in {90/2, 18/3, -54/4, -126/5, 162/6}{ \node[gv] (n\l) at (\a:1.4) {\l}; }
    \foreach \l in {2,...,6}{ \draw[gcut] (o) -- (n\l); }
    \foreach \x/\y in {2/3, 3/4, 4/5, 5/6, 6/2, 2/4, 3/5, 4/6, 5/2, 6/3}{ \draw[gfaded] (n\x) -- (n\y); }
    \node[note, anchor=west, text width=5.2cm, align=left] at (2.2,0) {$n = 6$: the $n-1 = 5$ red dashed bridges are exactly $\deg(1)$. With $k = 5$ the answer is 1; with $k = 4$ some bridge at island 1 survives and every island stays reachable: 6.};
\end{tikzpicture}
\end{center}

\begin{lstlisting}[style=cpp]
#include <bits/stdc++.h>
using namespace std;

int main() {
    int t;
    cin >> t;
    while (t--) {
        long long n, k;
        cin >> n >> k;
        cout << (k >= n - 1 ? 1 : n) << "\n";
    }
}
\end{lstlisting}
Official sample (pairs $(n,k)$): $(2,0),(2,1),(4,1),(5,10),(5,3),(4,4) \to 2, 1, 4, 1, 5, 1$. Verified also by trying every set of at most $k$ destroyed bridges for all $n \le 5$.
\end{solutionbox}

\begin{solutionbox}[E1.3 Maximum Total Importance of Roads]
Each road adds the values of its two endpoints, so city $v$ contributes $\text{val}(v)$ once per road: the total is $\sum_v \text{val}(v)\cdot\deg(v)$ (Tool 4). By the rearrangement inequality this is maximal when the values sorted increasingly meet the degrees sorted increasingly: the city with the $i$-th smallest degree gets value $i$. (Exchange argument: if $\deg(a) < \deg(b)$ but $\text{val}(a) > \text{val}(b)$, swapping the two values changes the total by $(\text{val}(a)-\text{val}(b))(\deg(b)-\deg(a)) > 0$.)
\begin{lstlisting}[style=cpp]
#include <bits/stdc++.h>
using namespace std;

class Solution {
public:
    long long maximumImportance(int n, vector<vector<int>>& roads) {
        vector<long long> deg(n, 0);
        for (auto& r : roads) { deg[r[0]]++; deg[r[1]]++; }
        sort(deg.begin(), deg.end());                   // smallest degree first
        long long total = 0;
        for (int i = 0; i < n; i++) total += deg[i] * (i + 1);   // value i+1
        return total;
    }
};
\end{lstlisting}
Example 1: degrees $2,3,4,2,1$; sorted $1,2,2,3,4$; total $1\cdot1 + 2\cdot2 + 2\cdot3 + 3\cdot4 + 4\cdot5 = 43$. Example 2: $\to 20$. The answer reaches about $5\cdot10^4 \cdot 2 \cdot 5\cdot10^4 = 5\cdot10^9$: \texttt{long long} is required.
\end{solutionbox}

\begin{solutionbox}[E1.4 Maximum Star Sum of a Graph]
A star is local to its center, so treat every vertex as the center. Adding a neighbour with value $\le 0$ never helps, so collect only \emph{positive} neighbour values, sort them decreasingly and take at most $k$. The star with zero edges (the center alone) is always allowed --- that is why the answer of Example 2 is $-5$.
\begin{lstlisting}[style=cpp]
#include <bits/stdc++.h>
using namespace std;

class Solution {
public:
    int maxStarSum(vector<int>& vals, vector<vector<int>>& edges, int k) {
        int n = vals.size();
        vector<vector<int>> good(n);                 // positive neighbour values per vertex
        for (auto& e : edges) {
            if (vals[e[1]] > 0) good[e[0]].push_back(vals[e[1]]);
            if (vals[e[0]] > 0) good[e[1]].push_back(vals[e[0]]);
        }
        int best = INT_MIN;
        for (int c = 0; c < n; c++) {
            sort(good[c].rbegin(), good[c].rend());  // largest first
            int sum = vals[c];
            for (int i = 0; i < (int)good[c].size() && i < k; i++) sum += good[c][i];
            best = max(best, sum);
        }
        return best;
    }
};
\end{lstlisting}
Example 1: center 3 (value 4) has neighbours with values $2, 10, -10, -20$; the two best positive ones give $4 + 10 + 2 = 16$. Total work $\bigO{|V| + |E|\log|E|}$; \texttt{nth\_element} instead of sorting gives $\bigO{|V|+|E|}$. Checked against trying every subset of neighbours on random small graphs.
\end{solutionbox}

\begin{solutionbox}[E1.5 Is it a degree sequence?]
\begin{itemize}[leftmargin=*]
    \item[(a)] $(3,2,2,2)$: the sum $9$ is odd --- impossible by the handshake lemma.
    \item[(b)] $(3,3,1,1)$: the sum $8$ is even, yet impossible. With 4 vertices, degree 3 means adjacent to \emph{all} others; two such vertices give each remaining vertex degree $\ge 2$.
    \item[(c)] $(3,3,2,2)$: yes --- $K_4$ minus one edge.
    \item[(d)] $(4,1,1,1,1)$: yes --- a star with 4 leaves.
\end{itemize}
\begin{center}
\begin{tikzpicture}[x=1cm, y=1cm]
    \begin{scope}
        \node[gv] (a) at (0,0) {3}; \node[gv] (b) at (1.5,0) {3}; \node[gv] (c) at (1.5,1.3) {2}; \node[gv] (d) at (0,1.3) {2};
        \draw[ge] (a) -- (b); \draw[ge] (a) -- (c); \draw[ge] (a) -- (d); \draw[ge] (b) -- (c); \draw[ge] (b) -- (d);
        \draw[gfaded, dashed] (c) -- (d);
        \node[font=\scriptsize\bfseries] at (0.75,-0.75) {(c) $K_4$ minus edge};
    \end{scope}
    \begin{scope}[xshift=4.5cm, yshift=0.65cm]
        \node[gv] (o) at (0,0) {4};
        \foreach \a in {45,135,225,315}{ \node[gv] (l\a) at (\a:1.1) {1}; \draw[ge] (o) -- (l\a); }
        \node[font=\scriptsize\bfseries] at (0,-1.4) {(d) star};
    \end{scope}
\end{tikzpicture}
\end{center}
(Labels are degrees.) An even sum is necessary but not sufficient; the exact test for simple graphs is the Erd\H{o}s--Gallai theorem, which we do not need here. All four answers were confirmed by enumerating every simple graph on 4 and 5 vertices.
\end{solutionbox}

\begin{solutionbox}[E1.6 Counting short paths]
(a) Fix the middle vertex $v$; the path is determined by an unordered pair of distinct edges at $v$: $\sum_v \binom{\deg(v)}{2}$. On the running example of the theory (degrees $2,3,2,2,1,0$) this is $1+3+1+1+0+0 = 6$.

(b) Fix the middle vertex $b$; choose an arrow into $b$ and an arrow out of $b$: $\sum_b \indeg{b}\cdot\outdeg{b}$. The two ends differ because no opposite pair $a \to b \to a$ exists. On the directed example of the theory (in/out: $0/2,\ 1/1,\ 3/1,\ 0/1,\ 1/0$): $0 + 1 + 3 + 0 + 0 = 4$, namely $0{\to}1{\to}2$, $0{\to}2{\to}4$, $1{\to}2{\to}4$, $3{\to}2{\to}4$.

Counting \emph{triangles} is a different story: the three edges are not all at one vertex, so degrees alone cannot do it (Chapter~22).
\end{solutionbox}

\begin{solutionbox}[E1.7 A tree with prescribed degrees]
Split the vertices into leaves ($d = 1$) and inner vertices ($d \ge 2$). For every inner vertex $v$: attach $d_v - 1$ leaves to it (each attached leaf is finished and never used again), after which $v$ needs exactly one more edge --- it has become a leaf itself. At the end exactly two leaves remain; join them.
\begin{lstlisting}[style=cpp]
#include <bits/stdc++.h>
using namespace std;

vector<pair<int,int>> treeFromDegrees(vector<int> d) {    // vertices 0..n-1
    int n = d.size();
    vector<pair<int,int>> edges;
    vector<int> leaves, inner;
    for (int v = 0; v < n; v++) (d[v] == 1 ? leaves : inner).push_back(v);
    for (int v : inner) {
        while (d[v] > 1) {                                  // attach a finished leaf to v
            int l = leaves.back(); leaves.pop_back();
            edges.push_back({l, v});
            d[v]--;
        }
        leaves.push_back(v);                                // v needs exactly one more edge
    }
    edges.push_back({leaves[0], leaves[1]});                // the last two leaves
    return edges;
}

int main() {
    vector<int> d = {1, 3, 1, 2, 1, 1, 3};                  // n = 7, sum = 12 = 2*7 - 2
    // prints (5,1) (4,1) (1,3) (3,6) (2,6) (0,6)
    for (auto& e : treeFromDegrees(d)) cout << "(" << e.first << "," << e.second << ") ";
    cout << "\n";
}
\end{lstlisting}
\begin{center}
\begin{tikzpicture}[x=1cm, y=1cm]
    \node[gv] (v1) at (0,0) {1}; \node[gv] (v3) at (1.6,0) {3}; \node[gv] (v6) at (3.2,0) {6};
    \node[gv] (v5) at (-1.1,0.8) {5}; \node[gv] (v4) at (-1.1,-0.8) {4};
    \node[gv] (v2) at (4.3,0.8) {2}; \node[gv] (v0) at (4.3,-0.8) {0};
    \draw[ge] (v5) -- (v1); \draw[ge] (v4) -- (v1); \draw[ge] (v1) -- (v3); \draw[ge] (v3) -- (v6);
    \draw[ge] (v2) -- (v6); \draw[gpath] (v0) -- (v6);
    \node[gdist, below=1mm of v1] {3}; \node[gdist, below=1mm of v3] {2}; \node[gdist, below=1mm of v6] {3};
    \node[note, anchor=west, text width=4.5cm, align=left] at (5.2,0) {degrees $1,3,1,2,1,1,3$ for vertices $0..6$; the orange edge joins the last two leaves 0 and 6.};
\end{tikzpicture}
\end{center}
\textbf{Why it never runs out of leaves.} $\sum_v (d_v - 1) = n - 2$, so the inner vertices need $\sum_{\text{inner}} (d_v - 1) = n - 2$ attachments in total, while the supply is the $L = n - |\text{inner}|$ original leaves plus one new leaf per finished inner vertex. Before inner vertex number $j$ the supply exceeds the demand because every later inner vertex still needs at least one attachment; at the end $L + |\text{inner}| - (n - 2) = 2$ leaves remain. Every edge removes a finished vertex for good, so no cycle can close, and $n - 1$ acyclic edges on $n$ vertices form a tree. Verified on every valid degree sequence with $n \le 7$.
\end{solutionbox}
